% Generated draft: requires mathematical and source review.
\documentclass[11pt]{article}
\usepackage[margin=25mm]{geometry}
\usepackage{fontspec}
\setmainfont{Latin Modern Roman}
\usepackage{amsmath,amssymb,mathtools,unicode-math}
\setmathfont{Latin Modern Math}
\usepackage{enumitem,needspace,xurl,hyperref}
\setlength{\parindent}{0pt}
\setlength{\parskip}{5pt}
\setlength{\emergencystretch}{3em}
\newcommand{\N}{\mathbb{N}}
\newcommand{\Z}{\mathbb{Z}}
\newcommand{\Q}{\mathbb{Q}}
\newcommand{\R}{\mathbb{R}}
\newcommand{\C}{\mathbb{C}}
\newcommand{\F}{\mathbb{F}}
\newcommand{\A}{\mathbb{A}}
\newcommand{\PP}{\mathbb{P}}
\newcommand{\E}{\mathbb{E}}
\newcommand{\eps}{\varepsilon}
\newcommand{\dd}{\,\mathrm{d}}
\begin{document}
\section*{UnsolvedMath ID: 30004033}
\textbf{Source problem code:} catalogue entry
\subsection{Definitions and mathematical statement}
The supplied question is formalized below under the conventional interpretation of graphs as finite, undirected, and simple. The source does not explicitly specify these conventions; extension to other graph domains is unresolved in this draft.

Let \(G=(V,E)\) be such a graph. The graph is planar if it admits a drawing in the plane in which edges meet only at common endpoints. It is triangle-free if it contains no cycle of length three. For \(v\in V\), let \(\deg_G(v)\) denote the number of vertices adjacent to \(v\). Let
\[
\mathcal I(G)=\{I\subseteq V: \text{no two distinct vertices of }I\text{ are adjacent}\}
\]
be the family of independent sets of \(G\). A fractional coloring assigns a weight \(w_I\in\R_{\geq 0}\) to every \(I\in\mathcal I(G)\), subject to
\[
\sum_{\substack{I\in\mathcal I(G)\\v\in I}} w_I\geq 1
\qquad\text{for every }v\in V.
\]
The fractional chromatic number is
\[
\chi_f(G)=\min\left\{\sum_{I\in\mathcal I(G)}w_I:
 w_I\geq 0\text{ for every }I\in\mathcal I(G),\quad
 \sum_{\substack{I\in\mathcal I(G)\\v\in I}}w_I\geq 1\text{ for every }v\in V\right\}.
\]
This definition gives \(\chi_f(G)=0\) when \(V=\varnothing\).

\textbf{Question.} Is it true that, for every finite simple undirected graph \(G=(V,E)\),
\[
\left(G\text{ is planar and triangle-free, and }
\deg_G(v)\leq 3\text{ for every }v\in V\right)
\quad\Longrightarrow\quad
\chi_f(G)\leq\frac{8}{3}\,?
\]
No connectedness, regularity, or lower bound on the number of vertices is required in this formulation.

The supplied record labels the problem partially solved and describes improved fractional upper bounds while retaining the proposed \(8/3\) bound as conjectural. These are supplied dated literature assessments, not independently verified mathematical or status claims.
\subsection{Short English statement}
Does every triangle-free planar graph in which each vertex has at most three neighbors have fractional chromatic number at most \(8/3\)? The displayed formulation uses finite simple undirected graphs, a convention not explicitly specified in the supplied source. The supplied partially solved label does not establish the proposed bound.
\subsection{Sources}
\begin{itemize}
\item Supplied local problem registry: permanent UnsolvedMath ID \(30004033\), title ``Fractional Coloring of Subcubic Triangle-Free Planar Graphs,'' Oberwolfach Reports Open Problems collection, exact source problem code \texttt{catalogue entry}. The record marks the statement as exact and describes it as retained byte-for-byte. No individual proposer is supplied; the proposed year is recorded as 2019.
\item Supplied report citation: \emph{Graph Theory (2020)}, Oberwolfach Reports, DOI 10.4171/owr/2019/1, pp.\ 5--63. \url{https://doi.org/10.4171/owr/2019/1}. The background describes Oberwolfach Reports as documenting mathematical workshops and their open-problem discussions.
\item Supplied literature reference: ``Heckman and Thomas, fractional chromatic number.'' \url{https://doi.org/10.1016/j.ejc.2013.06.006}.
\item Supplied literature reference: ``Large independent sets in triangle-free cubic graphs.'' \url{https://doi.org/10.19086/aic.13667}.
\item Supplied literature reference: ``11 over 4 colorability of subcubic triangle-free graphs.'' \url{https://www.advancesincombinatorics.com/article/133707-11-4-colorability-of-subcubic-triangle-free-graphs}.
\item Literature assessment dated \texttt{2026-08-22}: the supplied curation labels the problem partially solved and states that the fractional extension is incomplete, that later work improves fractional upper bounds, and that the natural \(8/3\)-type analogue remains conjectural. This is supplied dated curation derived from the OWR record, not an independently verified assessment. The supplied background explicitly requests independent verification of time-sensitive status and cited claims.
\end{itemize}
\end{document}
